\naslov{Greenove funkcije}

Naj bo $L$ linearen PDO s konstantnimi koeficienti na $\R^n$,
\[
  L = \sum_{\abs{\alpha} \le N} a_{\alpha} \partial^\alpha.
\]
Vzemimo $f \in \mathcal{C}^\infty_C$, in rešimo $L u = f$.

\begin{definicija}
  Funkcija $w$ se imenuje \pojem{fundamentalna rešitev} za operator $L$, če
  velja $Lw = \delta_0$, kjer je $\delta_0$ Diracova distribucija, definirana
  kot $\delta_0(f) = f(0)$.
\end{definicija}

Če je $w$ fundamentalna rešitev, je $w * f$ rešitev enačbe $Lu = f$; ker je $L$
linearen, velja
\[
  L(w * f) = (Lw)*f = \delta_0 * f = f,
\]
ker
\[
  \delta_0 * f(x) = \int \delta_0(x-t) f(t) dt = f(x).
\]

\vprasanje{Kaj je fundamentalna rešitev linearne PDE $Lu = f$? Kako z njo
  poiščeš rešitev $u$?}

Za fundamentalno rešitev prvo rešimo $Lw = 0$.

\begin{primer}
  Za $L = \Delta$ je operator invarianten na ortogonalne transformacije, torej
  lahko predpostavimo, da je $w(x) = \psi(\abs{x-\xi}^2)$ za nek fiksen $\xi \in
  \R^n$ ter funkcijo $\psi$.
  Enačbo lahko potem prevedemo na obliko
  \[
	\varphi''(r) + \frac{n-1}{r} \varphi'(r) = 0,
  \]
  kar lahko reduciramo na NDE prvega reda z rešitvijo $\varphi(r) = D + C r^{2 -
  n}$ za $n \ne 2$ oziroma $\varphi(r) = D + C \ln r$ za $n = 2$.
\end{primer}

Poglejmo si situacijo $n=2$ iz zgornjega primera.
Naj bo $M \subseteq \R^2$ omejena, odprta in povezana množica z gladkim robom.
Definiramo lahko funkcijo
\[
  \phi(z) = \inv{2 \pi} \log \abs{z}.
\]
Preverimo lahko, da je harmonična in lokalno integrabilna.
Definiramo še $k(z, \zeta) = \phi(z-\zeta)$ ter, za $f \in \zvezne{\cl{M}}$,
\[
  \Gamma(f)(z) = \iint_M f(\zeta) k(z, \zeta) dS
  = \iint_M f(\zeta) \phi(z-\zeta) dS
\]
Funkciji $\Gamma(f): M \to \R$ pravimo \pojem{Newtonov potencial}.

\begin{izrek}[Greenova reprezentacijska formula]
  Če je $u \in \zvodv{2}{\cl{M}}$ in $z \in M$, tedaj je
  \[
	u(z) = \int_{\partial M} \left( u \partial_{\vec{n}} k(z, \cdot) -
	  \partial_{\vec{n}} u \cdot k(z, \cdot) \right) ds
	+ \iint_M k(z, \cdot) \Delta u dS
  \]
\end{izrek}

\begin{proof}
  Skica.
  Za $\varepsilon > 0$ vzamemo $M_\varepsilon = M \brez \cl{K(z, \varepsilon)}$,
  ter uporabimo Greenovo identiteto
  \[
	\int_{\partial \Omega} (v \partial_{\vec{n}} u - u \partial_{\vec{n}} v) ds
	= \iint_\Omega (v \Delta u - u \Delta v) dS.
  \]
  Ostane integral leve strani enakosti iz izreka, po notranjem robu, za katerega
  pokažemo, da konvergira k $0$ za $\varepsilon \to 0$.
\end{proof}

\vprasanje{Povej Greenovo reprezentacijsko formulo.}

\begin{posledica}
  Za $u \in \mathcal{C}^2_C(\C)$ je $u = \Gamma(\Delta u)$.
\end{posledica}

Torej je $\phi$ fundamentalna rešitev za Laplaceov operator na $\R^2 \brez
\{0\}$.
Izrek bi želeli uporabiti za reševanje Cauchyjeve naloge $\Delta u = f$, $u = g$
na $\partial M$.
Problem je, da ne poznamo $\partial_{\vec{n}} u$ na $\partial M$, kar
potrebujemo v formuli.
Vpeljemo \pojem{korekcijski faktor}, namesto $k(z, \cdot)$ vzamemo $k(z, \cdot)
- h$ za primeren $h$.

Če v dokazu izreka namesto $k(z, \cdot)$ vzamemo neko funkcijo $H$ s $\Delta H =
0$, bi isti dokaz pokazal
\[
  \int_{\partial M} u \partial_{\vec{n}} H - H \partial_{\vec{n}} u
  = \iint_M H \Delta u,
\]
torej formula velja tudi, če $k$ nadomestimo s $k-H$.
Potem definiramo $G_z$ kot rešitev homogenega Dirichletovega problema $\Delta
G_z = 0$ in $G_z = \phi(\cdot - z)$ na $\partial M$.
Za tako izbiro bo $k(z, \cdot) - G_z = 0$ na $\partial M$, kar je točno to, kar
smo želeli od $h$.

\begin{definicija}
  Funkcija $G$, definirana kot
  \[
	G(z,w) = k(w,z) - G_z(w),
  \]
  se imenuje \pojem{Greenova funkcija} domene $M$ za Laplaceov operator.
\end{definicija}

\vprasanje{Izpelji Greenovo funkcijo za Laplaceov operator.}

\begin{izrek}
  Naj bo $G$ Greenova funkcija območja $M$, $z \in M$ ter $u \in
  \zvodv{2}{\cl{M}}$.
  Potem
  \[
	u(z) = \int_{\partial M} u(\zeta) \partial_{\vec{n}} G(z, \zeta) dS
	+ \iint_M G(z, \zeta) \Delta u(\zeta) dS.
  \]
\end{izrek}

% LocalWords:  PDO reprezentacijska reprezentacijsko
